The objective of this study was to estimate monthly depth-resolved compaction at the Tuku monitoring site during a six-month period in which finalised MLCW measurements were temporarily unavailable. One independent Bayesian ridge regression model was developed for each of six standardised 50~m depth sections using contemporaneous hydraulic head changes, vertical surface displacement increments, sediment composition, and seasonal terms as predictors.

The walk-forward evaluation showed that monthly compaction estimates reproduced the observed section-level MLCW increments with $R^2$ values ranging from \placeholder{MIN R2} to \placeholder{MAX R2} in sections S1 through S4. Performance was weaker in the two deepest sections, where the available piezometric observations did not screen the principal compacting deposits. Bayesian predictive intervals attained empirical coverage of \placeholder{COVERAGE RANGE} across the evaluated sections, providing calibrated uncertainty bounds for each monthly estimate from the first evaluation block onward. Under hypothetical reduced-frequency measurement schedules, the model distributed compaction among unobserved months with a monthly MAE of \placeholder{BEST H6 MAE} to \placeholder{WORST H6 MAE}~mm/month at the six-month interval, while cumulative endpoint errors remained within \placeholder{ENDPOINT RANGE}~mm of the observed field measurements.

These results demonstrated that contemporaneous hydraulic and surface displacement records can bridge a known data delay and provide provisional monthly estimates of section-level compaction at a single well-instrumented site. The model preserved direct MLCW observations as the reference measurement, using them to evaluate, constrain, and update the estimates when the finalised records became available. Extension to additional MLCW stations is required before the approach can support broader monitoring applications across the Choushui River Alluvial Fan.

\FloatBarrier
